\documentclass[10pt,a4paper]{amsart}
\usepackage{amsmath,amssymb,amsthm,graphicx}
\usepackage{verbatim}
\newenvironment{reference}{\comment}{\endcomment}
%\newenvironment{reference}{}{}
\newenvironment{slogan}{\comment}{\endcomment}
\newenvironment{history}{\comment}{\endcomment}

% For commutative diagrams we use Xy-pic
\usepackage[all]{xy}

% We use 2cell for 2-commutative diagrams.
\xyoption{2cell}
% Theorem environments.
%
\theoremstyle{plain}
\newtheorem{theorem}[subsection]{Theorem}
\newtheorem{proposition}[subsection]{Proposition}
\newtheorem{lemma}[subsection]{Lemma}

\theoremstyle{definition}
\newtheorem{definition}[subsection]{Definition}
\newtheorem{example}[subsection]{Example}
\newtheorem{exercise}[subsection]{Exercise}
\newtheorem{situation}[subsection]{Situation}

\theoremstyle{remark}
\newtheorem{remark}[subsection]{Remark}
\newtheorem{remarks}[subsection]{Remarks}

\numberwithin{equation}{subsection}

% Macros
%
\def\lim{\mathop{\mathrm{lim}}\nolimits}
\def\colim{\mathop{\mathrm{colim}}\nolimits}
\def\Spec{\mathop{\mathrm{Spec}}}
\def\Hom{\mathop{\mathrm{Hom}}\nolimits}
\def\Ext{\mathop{\mathrm{Ext}}\nolimits}
\def\SheafHom{\mathop{\mathcal{H}\!\mathit{om}}\nolimits}
\def\SheafExt{\mathop{\mathcal{E}\!\mathit{xt}}\nolimits}
\def\Sch{\mathit{Sch}}
\def\Mor{\mathop{\mathrm{Mor}}\nolimits}
\def\Ob{\mathop{\mathrm{Ob}}\nolimits}
\def\Sh{\mathop{\mathit{Sh}}\nolimits}
\def\NL{\mathop{N\!L}\nolimits}
\def\CH{\mathop{\mathrm{CH}}\nolimits}
\def\proetale{{pro\text{-}\acute{e}tale}}
\def\etale{{\acute{e}tale}}
\def\QCoh{\mathit{QCoh}}
\def\Ker{\mathop{\mathrm{Ker}}}
\def\Im{\mathop{\mathrm{Im}}}
\def\Coker{\mathop{\mathrm{Coker}}}
\def\Coim{\mathop{\mathrm{Coim}}}

% Boxtimes
%
\DeclareMathSymbol{\boxtimes}{\mathbin}{AMSa}{"02}

%
% Macros for moduli stacks/spaces
%
\def\QCohstack{\mathcal{QC}\!\mathit{oh}}
\def\Cohstack{\mathcal{C}\!\mathit{oh}}
\def\Spacesstack{\mathcal{S}\!\mathit{paces}}
\def\Quotfunctor{\mathrm{Quot}}
\def\Hilbfunctor{\mathrm{Hilb}}
\def\Curvesstack{\mathcal{C}\!\mathit{urves}}
\def\Polarizedstack{\mathcal{P}\!\mathit{olarized}}
\def\Complexesstack{\mathcal{C}\!\mathit{omplexes}}
% \Pic is the operator that assigns to X its picard group, usage \Pic(X)
% \Picardstack_{X/B} denotes the Picard stack of X over B
% \Picardfunctor_{X/B} denotes the Picard functor of X over B
\def\Pic{\mathop{\mathrm{Pic}}\nolimits}
\def\Picardstack{\mathcal{P}\!\mathit{ic}}
\def\Picardfunctor{\mathrm{Pic}}
\def\Deformationcategory{\mathcal{D}\!\mathit{ef}}
\title{A delta-functor vanishing on suitable embeddings is universal}
\author{The Stacks Project Authors}
\date{stacks-project:homology.tex:local-snapshot-20261012; GNU FDL 1.2 or later}
\begin{document}
\maketitle
\noindent\textit{Editorial note.} All delivered mathematical text is the author's own, copied from the pinned Stacks Project chapter file homology.tex; no statement or proof step was rewritten, reordered or removed.\par\medskip
\noindent\textit{Editorial note.} The proof is the author's complete argument; it ends with 'We omit the verification that this defines a transformation of functors.', which is the routine functoriality check on components already constructed in the span (LANE\_BRIEF 5, routine closing verification => collect), and the sentence is preserved verbatim. The corpus counts comparable proofs: STK-algebra-adjoint-tensor counts sources/stacks-algebra lines 2695-2702 containing 'We omit the verification that these constructions are mutually inverse.', and STK-algebra-loc-and-modules counts lines 1345-1353 of the same file.\par\medskip
\noindent\textit{Editorial note.} Every labelled reference occurring inside the counted spans is carried in assembly.external\_references as a non-load-bearing navigation pointer (same-chapter ones at their definition site, cross-chapter ones at the invocation inside the counted span), so no load-bearing mathematics is replaced by a stub. No bibliography citation occurs inside the counted spans.\par\medskip
\noindent\textit{Editorial note.} The same-name check for this card scanned every staged Stacks chapter of this repository -- all sources/stacks-*/pinned-original-source.tex files, the card's own chapter excluded -- and the labels found are recorded in dedup.same\_name\_cross\_chapter on the card.\par\medskip
\noindent\textit{Editorial note.} The preamble is the chapter's own native preamble excerpt already used for this source: the comment-type reference, slogan and history environments (native preamble lines 8-13), the Xy-pic setup (native preamble lines 14-18) and the native theorem declarations with the mathematical macros (native preamble lines 154-217, the actual last line of the file).\par\medskip
\noindent\textit{Allowed standard prerequisites (editorial).}\begin{itemize}
\item Delta-functors and the universal property. Same chapter, the definitions of a delta-functor and of a universal delta-functor. Supplies the notion under discussion and the data, a transformation of functors from the degree zero part, that the construction has to extend uniquely.
\item The long exact cohomology sequence of a delta-functor. Same chapter, the exactness properties of a delta-functor. Supplies the identification of the next degree of the functor with the cokernel of the map induced by the chosen embedding, which is what makes the induction step work.
\end{itemize}\medskip
\begin{lemma}
\label{lemma-efface-implies-universal}
Let $\mathcal{A}, \mathcal{B}$ be abelian categories.
Let $F = (F^n, \delta_F)$ be a $\delta$-functor
from $\mathcal{A}$ to $\mathcal{B}$.
Suppose that for every $n > 0$ and any $A \in \Ob(\mathcal{A})$
there exists an injective morphism $u : A \to B$ (depending on $A$ and $n$)
such that $F^n(u) : F^n(A) \to F^n(B)$ is zero. Then $F$ is a universal
$\delta$-functor.
\end{lemma}

\begin{proof}
Let $G = (G^n, \delta_G)$ be a $\delta$-functor
from $\mathcal{A}$ to $\mathcal{B}$ and let $t : F^0 \to G^0$
be a morphism of functors. We have to show there exists
a unique morphism of $\delta$-functors $\{t^n\}_{n \geq 0} : F \to G$
such that $t = t^0$. We construct $t^n$ by induction on $n$.
For $n = 0$ we set $t^0 = t$.
Suppose we have already constructed a unique sequence of
transformation of functors $t^i$ for $i \leq n$ compatible with
the maps $\delta$ in degrees $\leq n$.

\medskip\noindent
Let $A \in \Ob(\mathcal{A})$. By assumption we may choose
an embedding $u : A \to B$ such that $F^{n + 1}(u) = 0$.
Let $C = B/u(A)$. The long exact cohomology sequence for
the short exact sequence $0 \to A \to B \to C \to 0$ and the
$\delta$-functor $F$ gives that
$F^{n + 1}(A) = \Coker(F^n(B) \to F^n(C))$ by our choice of $u$.
Since we have already defined $t^n$ we can set
$$
t^{n + 1}_A : F^{n + 1}(A) \to G^{n + 1}(A)
$$
equal to the unique map such that
$$
\xymatrix{
\Coker(F^n(B) \to F^n(C)) \ar[r]_{t^n}
\ar[d]_{\delta_{F, A \to B \to C}} &
\Coker(G^n(B) \to G^n(C))
\ar[d]^{\delta_{G, A \to B \to C}} \\
F^{n + 1}(A) \ar[r]^{t^{n + 1}_A} &
G^{n + 1}(A)
}
$$
commutes. This is clearly uniquely determined by the requirements
imposed. We omit the verification that this defines a transformation
of functors.
\end{proof}

\end{document}
